% !TEX root = cap05_standalone.tex
% Capítulo 5. Modelo de crecimiento urbano
% Estructura según índice de la Dra. Lárraga.

\chapter{Modelo de crecimiento urbano}
\label{ch:modelo-crecimiento-urbano}

% ---- Contenido migrado (copy/paste) desde report/tesis/book/cap04/cap04.tex ----

\begin{resumen}
Este capítulo muestra una visión exhaustiva del modelo citado a lo largo de todo este trabajo como WoE-A, se presentan las decisiones explícitas de cada componente del modelo, preprocesameinto y validación. Si bien el segundo capítulo del presente trabajo da una definición a cada concepto utilizado, no hay una argumentación de por qué se utiliza cada una de las componentes del flujo completo que al trabajo presente le compete. Se inicia, como ya el lector debe saber, con la definición del flujo de preprocesamiento de datos, seguido por el modelo que consume los datos preprocesados para calcular la matriz de transiciones históricas, seguido por la definición del autómata celular probabilístico y finalmente la validación del modelo.
\end{resumen}

\section{Proprocesamiento de datos.}

El preprocesamiento de datos, para este trabajo sienta la verdad que el modelo WoE-AC utiliza para predecir la dinámica urban. Por ende, se debe entender el preprocesamiento de datos como el proceso que convierte las capturas RGB en mapas binarios urbano/no-urbano (pipeline LBP\,+\,K-Means\,+\,SVM); en el capítulo \ref{ch:marco-teorico} en particular en la figura \ref{fig:pipeline-rgb-to-binary} se muestra el flujo completo de preprocesamiento de datos, primero: se obtiene un tensor RGB de $1{,}792 \times 3{,}024 \times 3$ píxeles, esto constituye la información cruda de entrada al flujo. 

En el siguiente paso, se realiza una transformacion de este arreglo de $1{,}792 \times 3{,}024 \times 3$ píxeles a un arreglo de $1{,}792 \times 3{,}024 \times 23$ píxeles. Este paso conlleva la primer decisión critica en la definición del flujo, pues en vez de utilizar índices espectrales, se decide obtener 23 features por celda. Esto tiene una reazón de ser: las capturas disponibles son RGB, de tres bandas visibles, NDVI requiere NIR; NDWI requiere NIR o SWIR; NDBI requiere SWIR. Ninguno de esos índices es computable a partir de RGB puro, eso elimina la clase mas común de clasificadores de cobertura de suelo que usa la literatura. Esto tiene una consecuencia directa, con solo bandas visibles, los píxeles urbanos (techo gris, asfalto) y los de suelo desnudo (arcilla, tierra agrícola seca) tienen un parecido espectral RGB prácticamente indistinguibles. La solución a este problema consiste en extender la representación por pixel con 23 features definidas en las tablas \ref{tab:features_rgb}, \ref{tab:features_textura} y \ref{tab:features_espectrales}. El objetivo es codificar la textura (LBP, Sobel, entropía, contraste) y espacios de color perceptuales (LAB, HSV). El tejido urbano tiene un patrón espacial regular, cuadrícula de calles, alineación de techados; que la vegetación y el suelo agrícola no tienen. LBP captura exactamente esa regularidad; es un descriptor diseñado para distinguir texturas y ha sido validado en clasificación de imágenes de satélite de baja resolución espectral. 

Sin embargo, al ser 23 features o características, esto conlleva un problema de explosión de parámetros a clasificar, es decir, en vez de trabajar sobre un espacio de $\mathbb{R}^{3}$, se debe trabajar sobre un espacio de $\mathbb{R}^{23}$, lo cual es computacionalmente muy demandante. La solución a este problema consiste en reducir la dimensionalidad de las features mediante PCA.

Antes de aplicar el PCA hay un paso intermedio que no puede omitirse: las 23 características viven en escalas muy distintas. Un canal RGB está acotado entre $0$ y $255$, mientras que el índice de exceso de verde o las componentes del espacio LAB toman valores en rangos completamente diferentes. Si se aplicara PCA directamente sobre esas escalas, las features de mayor magnitud dominarían artificialmente las primeras componentes, no porque contengan más información, sino simplemente por tener números más grandes. Por ello cada característica se estandariza primero con un \texttt{StandardScaler} (Algoritmo~\ref{alg:standard-scaler}), que centra cada columna en media cero y varianza unitaria; con esto las 23 features quedan en condiciones comparables y el PCA puede operar sobre la varianza real de cada una.

Sobre la matriz ya estandarizada se aplica entonces el PCA (Algoritmo~\ref{alg:pca}), que proyecta las 23 dimensiones a un espacio de 8 componentes ortogonales. La elección de 8 componentes no es arbitraria: durante el desarrollo del modelo se observó que estas 8 componentes retienen alrededor del $96\%$ de la varianza total de la serie. Este número revela además un hecho importante sobre los datos: si las features de textura y color fueran meras combinaciones del RGB crudo, bastarían 3 o 4 componentes para alcanzar ese mismo $96\%$; el que se necesiten 8 confirma que la textura y el color sí aportan información complementaria y no redundante. La ortogonalidad de las componentes resuelve, de paso, un problema que de otro modo afectaría al siguiente paso del flujo: la fuerte correlación entre R, G, B, intensidad y brillo distorsionaría la distancia euclidiana con la que K-Means agrupa, al contar varias veces la misma información espectral. Tras el PCA se obtiene una matriz $Z \in \mathbb{R}^{N \times 8}$ que es la entrada del clasificador.

La siguiente decisión es cómo etiquetar cada píxel como urbano o no-urbano. La vía más directa sería una clasificación supervisada con etiquetas manuales, pero aquí aparece una restricción propia del problema: la serie cubre 37 años (1984--2020). Etiquetar manualmente exigiría entre 500 y $1{,}000$ píxeles por imagen, coherentes entre años, lo que suma decenas de miles de decisiones manuales. Por eso se opta por una vía no supervisada: K-Means con $k=2$ (Algoritmo~\ref{alg:kmeans}). Su valor en este contexto no es la exactitud cartográfica frente a una verdad terreno, sino la consistencia interna de la serie, minimizar la inercia intra-cluster, es decir, es idéntica para todos los años, el criterio que separa lo urbano de lo no-urbano en 1990 es el mismo que en 2010. Esa consistencia es justamente lo que hace que las transiciones entre años sean comparables y puedan, más adelante, alimentar el cálculo WoE sin introducir sesgos por cambios de criterio.

K-Means por sí solo, sin embargo, deja una frontera de decisión frágil: al separar el espacio por la mediatriz entre los dos centroides, los píxeles de valor intermedio, zonas periurbanas, sombras, cobertura mixta, pueden desplazar esos centroides hacia regiones ambiguas. Para corregir esto, las etiquetas de K-Means no se usan como resultado final, sino como pseudo-etiquetas para entrenar un SVM lineal (Algoritmo~\ref{alg:svm}). El SVM busca el hiperplano de máximo margen en el espacio PCA, lo que regulariza la frontera y la vuelve robusta frente a esos píxeles atípicos; el kernel lineal es suficiente porque el PCA ya proyectó los datos en un espacio donde las dos clases son aproximadamente separables. El efecto de este refinamiento es medible: a lo largo de la serie el SVM alcanzó entre el $88\%$ y el $94\%$ de exactitud sobre el $20\%$ de píxeles reservados para prueba, por encima de lo que produce una umbralización simple sobre las mismas imágenes.

El último paso es operativo, pero necesario para cerrar el flujo: el SVM entrega dos clases, pero no sabe cuál de ellas corresponde a la mancha urbana. Esa asignación se realiza de forma manual, fijando la convención urbano${}=1$ / no-urbano${}=0$. Con ello, cada captura RGB queda convertida en un mapa binario $U_t(x,y) \in \{0,1\}$, y la repetición del procedimiento sobre los 37 años produce la serie $\{U_t\}_{t=1984}^{2020}$, que es el insumo del que parte todo el modelo WoE-AC: a partir de ella se calculan las transiciones históricas, se derivan las siete variables espaciales y se entrena el Weight of Evidence.

La elección de K-Means sobre otros algoritmos no supervisados, como los modelos de mezclas gaussianas (GMM), el agrupamiento jerárquico o el basado en densidad (DBSCAN), responde a la escala del problema: cada imagen contiene del orden de $5{,}4$ millones de píxeles, y métodos como el agrupamiento jerárquico o el espectral exigen matrices de afinidad de tamaño $N \times N$, inviables en memoria a esa escala. K-Means, en cambio, escala de forma lineal con el número de píxeles y, fijada la semilla y el número de inicializaciones, produce una partición estable y reproducible. GMM y DBSCAN añadirían supuestos sobre la forma o la densidad de los grupos que no se justifican aquí, donde lo que se busca no es descubrir estructura compleja sino una separación binaria estable y consistente a lo largo de los 37 años de la serie.

\section{Modelo de crecimiento urbano.}

